\documentclass[11pt]{amsart}

\usepackage{amsmath,amssymb,amsthm}
\usepackage{graphicx}
\usepackage{hyperref}
\usepackage{booktabs}
\usepackage{microtype}  % microtypography

\theoremstyle{plain}
\newtheorem{theorem}{Theorem}
\newtheorem{proposition}[theorem]{Proposition}
\newtheorem{lemma}[theorem]{Lemma}
\newtheorem{corollary}[theorem]{Corollary}

\theoremstyle{definition}
\newtheorem{definition}[theorem]{Definition}
\newtheorem{example}[theorem]{Example}

\theoremstyle{remark}
\newtheorem{remark}[theorem]{Remark}

\newcommand{\ee}{\mathrm{e}}
\newcommand{\ii}{\mathrm{i}}
\newcommand{\R}{\mathbb{R}}
\newcommand{\C}{\mathbb{C}}
\newcommand{\Z}{\mathbb{Z}}
\newcommand{\Kfrac}{\mathop{\mathrm{K}}}

\title{A Spiral for Euler's Number}
\author{J. Popelka}
\date{\today}

\begin{document}

\begin{abstract}
We present a monotonically convergent rational series for Euler's number~$\ee$
derived from classical continued fraction theory, together with its analytic
continuation via modified Bessel functions. The resulting formula exhibits an
elegant symmetry and traces a remarkable closed spiral in the complex plane.
While the connection between~$\ee$ and Bessel polynomials is classical, the
symmetric Bessel formula and its geometric visualization appear to be new.
\end{abstract}

\maketitle

%% ============================================================
\section{Introduction}
%% ============================================================

Consider the parametric curve in the complex plane defined by
\begin{equation}\label{eq:g}
g(t) = \frac{-16\pi\ee \cdot t}
       {K_{2t-1}\!\left(-\tfrac{1}{2}\right) \cdot K_{2t+1}\!\left(-\tfrac{1}{2}\right)},
\end{equation}
where $K_\nu(z)$ is the modified Bessel function of the second kind.

This curve, shown in Figure~\ref{fig:spiral}, traces a double spiral that closes
smoothly at the origin---a topological circle embedded in~$\C$. More remarkably,
sampling this curve at quarter-integer points recovers Euler's number:
\begin{equation}\label{eq:e-series}
\ee = 1 + \sum_{n=0}^{\infty} g\!\left(\tfrac{3}{4} + n\right).
\end{equation}

The series converges rapidly, gaining approximately six decimal digits per term.
In this note, we derive formula~\eqref{eq:g} from classical continued fraction
theory and explore its geometric properties.

\begin{figure}[ht]
\centering
% \includegraphics[width=0.6\textwidth]{figures/spiral.pdf}
\fbox{\parbox{0.6\textwidth}{\centering\vspace{3cm}
[Parametric plot of $g(t)$ for $t \in [-3, 3]$]
\vspace{3cm}}}
\caption{The spiral $t \mapsto g(t) \in \C$. The curve is odd-symmetric about
the origin and closes smoothly as $t \to \pm\infty$.}
\label{fig:spiral}
\end{figure}

%% ============================================================
\section{From Continued Fractions to a Monotonic Series}
%% ============================================================

\subsection{The arithmetic continued fraction}

Euler discovered that
\begin{equation}\label{eq:coth-cf}
\coth\!\left(\tfrac{1}{2}\right) = \frac{\ee + 1}{\ee - 1} = [2; 6, 10, 14, 18, \ldots],
\end{equation}
a continued fraction with arithmetic progression. Subtracting~2 gives the
purely periodic form
\begin{equation}\label{eq:cf-arith}
\frac{3 - \ee}{\ee - 1} = [0; 6, 10, 14, 18, \ldots]
= \Kfrac_{k=1}^{\infty} \frac{1}{4k+2}.
\end{equation}

Let $p_n/q_n$ denote the $n$-th convergent of~\eqref{eq:cf-arith}. The standard
recurrence gives
\begin{equation}\label{eq:recurrence}
p_n = (4n+2)\, p_{n-1} + p_{n-2}, \qquad
q_n = (4n+2)\, q_{n-1} + q_{n-2},
\end{equation}
with initial values $p_0 = 0$, $p_1 = 1$, $q_0 = 1$, $q_1 = 6$.

\subsection{The transformed convergents}

Define the \emph{transformed convergent}
\begin{equation}\label{eq:T-def}
T_n = \frac{3q_n + p_n}{q_n + p_n}.
\end{equation}

\begin{proposition}
$T_n$ equals the $(3n+2)$-th convergent of the standard continued fraction
for~$\ee$.
\end{proposition}

\begin{proof}
This follows from the M\"obius transformation
$\sigma(x) = (1-x)/(1+x)$ applied to $(3-\ee)/(\ee-1)$, which yields $\ee - 2$.
\end{proof}

The sequence $(T_n)$ alternates around~$\ee$: odd-indexed terms undershoot,
even-indexed terms overshoot. The first few values are:
\[
T_1 = \frac{19}{7}, \quad
T_2 = \frac{193}{71}, \quad
T_3 = \frac{2721}{1001}, \quad
T_4 = \frac{49171}{18089}, \quad \ldots
\]

\subsection{Telescoping to monotonicity}

The difference between consecutive transforms is
\begin{equation}\label{eq:diff}
T_{n+1} - T_n = \frac{(-1)^{n+1} \cdot 2}{s_n \cdot s_{n+1}},
\end{equation}
where $s_n = p_n + q_n$ satisfies the recurrence~\eqref{eq:recurrence} with
$s_0 = 1$, $s_1 = 7$.

Pairing consecutive differences eliminates the alternation:
\begin{align}
(T_{2k} - T_{2k-1}) + (T_{2k+1} - T_{2k})
&= \frac{2}{s_{2k-1} s_{2k}} - \frac{2}{s_{2k} s_{2k+1}} \notag\\
&= \frac{2(s_{2k+1} - s_{2k-1})}{s_{2k-1} \cdot s_{2k} \cdot s_{2k+1}}.
\label{eq:paired}
\end{align}

Using the recurrence, $s_{2k+1} - s_{2k-1} = (8k+6)\, s_{2k}$, so the middle
term cancels:
\begin{equation}\label{eq:simplified-term}
\frac{2(8k+6)}{s_{2k-1} \cdot s_{2k+1}} = \frac{4(4k+3)}{s_{2k-1} \cdot s_{2k+1}}.
\end{equation}

Re-indexing with $j = k - 1$ yields our main rational formula.

\begin{theorem}[Monotonic series for $\ee$]\label{thm:main}
\begin{equation}\label{eq:monotonic}
\ee = 1 + 4\sum_{j=0}^{\infty} \frac{4j+3}{s_{2j-1} \cdot s_{2j+1}},
\end{equation}
where $s_{-1} = 1$, $s_1 = 7$, $s_3 = 1001$, $s_5 = 398959$, \ldots{} are the
odd-indexed terms of the sequence satisfying
$s_n = (4n+2)\,s_{n-1} + s_{n-2}$.
\end{theorem}

Each term adds approximately six correct decimal digits:

\begin{center}
\begin{tabular}{ccc}
\toprule
Terms & Partial sum & Correct digits \\
\midrule
1 & $19/7 = 2.714\ldots$ & 2 \\
2 & $2721/1001 = 2.71828171\ldots$ & 6 \\
3 & $1084483/398959 = 2.718281828458\ldots$ & 12 \\
4 & $851568301/313191553 = 2.71828182845904523\ldots$ & 18 \\
\bottomrule
\end{tabular}
\end{center}

%% ============================================================
\section{The Bessel Connection}
%% ============================================================

\subsection{Bessel polynomials}

The sequence $(s_n)$ connects to \emph{Bessel polynomials} $y_n(x)$, defined by
\begin{equation}
y_n(x) = \sum_{k=0}^{n} \frac{(n+k)!}{(n-k)!\, k!} \left(\frac{x}{2}\right)^k
= {}_2F_0\!\left(-n, n+1; \,; -\frac{x}{2}\right).
\end{equation}

\begin{proposition}
$s_n = (-1)^{n+1} y_{n+1}(-2)$.
\end{proposition}

The values $y_n(-2)$ form OEIS sequence A002119:
\[
1, -1, 7, -71, 1001, -18089, 398959, \ldots
\]

\subsection{Modified Bessel functions}

For half-integer order $\nu = n + \frac{1}{2}$, the modified Bessel function
$K_\nu(z)$ has the elementary form~\cite{Watson1944}
\begin{equation}\label{eq:K-half}
K_{n+\frac{1}{2}}(z) = \sqrt{\frac{\pi}{2z}}\, \ee^{-z} \cdot P_n(1/z),
\end{equation}
where $P_n$ is a polynomial related to Bessel polynomials.

At the special argument $z = -\frac{1}{2}$ (on the branch cut), we obtain
\begin{equation}\label{eq:K-special}
K_{n+\frac{1}{2}}\!\left(-\tfrac{1}{2}\right) = \ii\sqrt{\pi\ee} \cdot y_n(-2).
\end{equation}

This is pure imaginary! Consequently, the \emph{product} of two such values
is real:
\begin{equation}
K_\mu\!\left(-\tfrac{1}{2}\right) \cdot K_\nu\!\left(-\tfrac{1}{2}\right)
= -\pi\ee \cdot y_{\mu-1/2}(-2) \cdot y_{\nu-1/2}(-2).
\end{equation}

%% ============================================================
\section{The Symmetric Formula}
%% ============================================================

Define the \emph{term function}
\begin{equation}\label{eq:term-j}
f(j) = \frac{4(4j+3)}{s_{2j-1} \cdot s_{2j+1}}, \qquad j \in \Z_{\geq 0}.
\end{equation}

Using~\eqref{eq:K-special}, this extends to a meromorphic function on~$\C$:
\begin{equation}\label{eq:f-bessel}
f(j) = \frac{-4(4j+3)\pi\ee}
       {K_{2j+\frac{1}{2}}\!\left(-\tfrac{1}{2}\right) \cdot
        K_{2j+\frac{5}{2}}\!\left(-\tfrac{1}{2}\right)}.
\end{equation}

\subsection{Canonical form}

The function $f(j)$ is odd about $j = -\frac{3}{4}$. Substituting
$t = j + \frac{3}{4}$ yields the \emph{canonical form}
\begin{equation}\label{eq:g-canonical}
\boxed{g(t) = \frac{-16\pi\ee \cdot t}
       {K_{2t-1}\!\left(-\tfrac{1}{2}\right) \cdot K_{2t+1}\!\left(-\tfrac{1}{2}\right)}}
\end{equation}
which is odd about the origin: $g(-t) = -g(t)$.

\begin{remark}
The Bessel orders $2t \pm 1$ are symmetric about~$2t$, giving the formula
its elegant structure. The original term function is recovered via
$f(j) = g(j + \frac{3}{4})$.
\end{remark}

\subsection{Equivalent forms}

Different symmetry centers yield equivalent formulas:

\begin{center}
\begin{tabular}{ccc}
\toprule
Center & Variable & Functional equation \\
\midrule
$j = -\frac{3}{4}$ & $f(j)$ & $f(-\frac{3}{4} + t) = -f(-\frac{3}{4} - t)$ \\[4pt]
$t = 0$ & $g(t)$ & $g(-t) = -g(t)$ \\[4pt]
$s = \frac{1}{2}$ & $h(s) = g(s - \frac{1}{2})$ & $h(1-s) = -h(s)$ \\
\bottomrule
\end{tabular}
\end{center}

The last form, with symmetry about $s = \frac{1}{2}$, is reminiscent of
functional equations in analytic number theory.

%% ============================================================
\section{The Spiral}
%% ============================================================

\subsection{Parametric curve}

The map $t \mapsto g(t) \in \C$ traces a double spiral (Figure~\ref{fig:spiral}).

\begin{proposition}[Properties of the spiral]
\leavevmode
\begin{enumerate}
\item \textbf{Odd symmetry:} $g(-t) = -g(t)$ (180° rotation about origin).
\item \textbf{Zero at origin:} $g(0) = 0$.
\item \textbf{Decay:} $\lim_{t \to \pm\infty} g(t) = 0$.
\end{enumerate}
\end{proposition}

\subsection{Topology: a smooth closed curve}

\begin{theorem}
The curve $\gamma: \R \cup \{\infty\} \to \C$ defined by $g$ is homeomorphic
to~$S^1$.
\end{theorem}

\begin{proof}
We verify smooth closure at both the origin and infinity.

\textbf{At infinity:} As $t \to +\infty$, using~\eqref{eq:K-half},
$g(t) \sim c \cdot t \cdot \ee^{-2|t|}$ for some constant~$c$.
The curve approaches the origin along the negative real axis (angle~$\pi$).
Similarly, $t \to -\infty$ approaches from the positive real axis (angle~$0$).
These directions are antipodal, so the curve closes smoothly at infinity.

\textbf{At the origin:} We have $g(0) = 0$ and
\[
g'(0) = \lim_{t \to 0} \frac{g(t)}{t} = \frac{-16\pi\ee}{K_{-1}(-\frac{1}{2}) \cdot K_1(-\frac{1}{2})},
\]
which is the same finite nonzero value from both sides.
\end{proof}

The spiral is thus a smooth embedding of the circle into the plane, with:
\begin{itemize}
\item The point at infinity mapping to the origin (from opposite directions)
\item The origin $t = 0$ mapping to $g(0) = 0$ with well-defined tangent
\end{itemize}

\subsection{Recovery of $\ee$}

The series~\eqref{eq:monotonic} samples the spiral at $t = \frac{3}{4} + n$:
\begin{equation}
\ee = 1 + \sum_{n=0}^{\infty} g\!\left(\tfrac{3}{4} + n\right).
\end{equation}

These sample points lie between consecutive real-axis crossings of the curve.

\subsection{When is $g(t)$ real?}

The function $g(t)$ takes real values precisely when both Bessel orders
$2t \pm 1$ are half-integers---that is, when the orders correspond to
\emph{spherical} Bessel functions.

\begin{proposition}
$g(t) \in \R$ if and only if $t \in \{\frac{2k+1}{4} : k \in \Z\} \cup \{0\}$.
\end{proposition}

The real-axis crossings occur at
$t = \ldots, -\frac{5}{4}, -\frac{3}{4}, -\frac{1}{4}, 0, \frac{1}{4}, \frac{3}{4}, \frac{5}{4}, \ldots$

%% ============================================================
\section{Discussion}
%% ============================================================

\subsection{What is known}

The connection between $\ee$ and Bessel polynomials dates to the 1950s;
see OEIS A002119~\cite{OEIS-A002119} and references therein. The continued
fraction~\eqref{eq:coth-cf} was discovered by Euler in~1737~\cite{Euler1737}.

The convergence rate of CF-based methods for computing~$\ee$ is well documented;
see~\cite{GourdonSebah} for a comprehensive treatment.

\subsection{What appears new}

\begin{enumerate}
\item The symmetric formula~\eqref{eq:g-canonical} with its $K_{2t\pm 1}$
structure.
\item The spiral visualization and its topological properties.
\item The explicit characterization of real-axis crossings via half-integer
orders.
\end{enumerate}

\subsection{Spherical interpretation}

Half-integer Bessel orders correspond to \emph{spherical} Bessel functions,
which arise in three-dimensional problems (multipole expansions, quantum
scattering). Integer orders correspond to \emph{cylindrical} geometry.

The fact that our formula uses half-integer orders suggests---at least
metaphorically---that $\ee$ ``lives'' in spherical rather than cylindrical
Bessel space. We offer this as a pedagogical observation rather than a
deep theorem.

%% ============================================================
\section{Conclusion}
%% ============================================================

Starting from Euler's classical continued fraction for $\coth(\frac{1}{2})$,
we derived a monotonically convergent rational series for~$\ee$ and found
its natural analytic continuation via modified Bessel functions. The
resulting formula~\eqref{eq:g-canonical} exhibits an elegant symmetry and
traces a closed spiral in the complex plane.

The spiral provides a geometric perspective on one of mathematics' most
fundamental constants: $\ee$ emerges from sampling a smooth closed curve
at quarter-integer points.

%% ============================================================
\bibliographystyle{amsplain}
\bibliography{references}

\end{document}
